% BEGIN VOL07-EXPANSION VII19
\section{Supplementary worked examples}
\label{sec:vii19-pedagogy}

\begin{example}[Planes are minimal]\label{ex:vii19-ped-01}
Every plane has constant unit normal and zero shape operator, so both principal curvatures vanish.  Hence \(H=0\).  Planes are the simplest minimal surfaces and the local model around which the minimal-surface equation linearizes.
\end{example}

\begin{example}[The catenoid]\label{ex:vii19-ped-02}
Rotating a catenary about its axis produces the catenoid, a nonplanar minimal surface of revolution.  Its two principal curvatures are equal in magnitude and opposite in sign, so \(H=0\) while \(K<0\).
\end{example}

\begin{example}[The helicoid]\label{ex:vii19-ped-03}
The helicoid is both ruled and minimal.  It provides a striking contrast with the catenoid: the two surfaces look very different extrinsically, yet both satisfy the same variational condition \(H=0\).  Locally they are connected through the classical associated family of minimal surfaces.
\end{example}

\section{Graded supplementary exercises}
The following exercises add a deliberate progression from concrete calculation to structural proof, hypothesis testing, application, and synthesis.

\subsection*{Standard computations and constructions}

\begin{exercise}[Plane mean curvature]\label{exr:vii19-ped-01}
Compute the mean curvature of a plane.
\end{exercise}
\begin{hint}
Its shape operator is zero.
\end{hint}
\begin{solution}
\(H=0\).
\end{solution}

\begin{exercise}[Opposite principal curvatures]\label{exr:vii19-ped-02}
If \(k_1=a\) and \(k_2=-a\), compute \(H\) and \(K\).
\end{exercise}
\begin{hint}
Average and multiply.
\end{hint}
\begin{solution}
\(H=0\) and \(K=-a^2\).
\end{solution}

\begin{exercise}[Cylinder minimality]\label{exr:vii19-ped-03}
Is a circular cylinder minimal?
\end{exercise}
\begin{hint}
Use principal curvatures \(1/R\) and \(0\).
\end{hint}
\begin{solution}
No.  Its mean curvature has nonzero magnitude \(1/(2R)\).
\end{solution}

\begin{exercise}[Minimal graph equation]\label{exr:vii19-ped-04}
State the minimal-surface equation for a graph \(z=f(x,y)\).
\end{exercise}
\begin{hint}
Use divergence form.
\end{hint}
\begin{solution}
\(\operatorname{div}\left(\nabla f/\sqrt{1+|\nabla f|^2}\right)=0\), equivalently \((1+f_y^2)f_{xx}-2f_xf_yf_{xy}+(1+f_x^2)f_{yy}=0\).
\end{solution}

\begin{exercise}[Linearization]\label{exr:vii19-ped-05}
What equation results by linearizing the minimal graph equation at \(f=0\)?
\end{exercise}
\begin{hint}
Discard terms quadratic in first derivatives.
\end{hint}
\begin{solution}
The Laplace equation \(f_{xx}+f_{yy}=0\).
\end{solution}

\subsection*{Proofs}

\begin{exercise}[Orientation independence of minimality]\label{exr:vii19-ped-06}
Prove the condition \(H=0\) does not depend on normal orientation.
\end{exercise}
\begin{hint}
Normal reversal changes \(H\) to \(-H\).
\end{hint}
\begin{solution}
Zero is unchanged by sign, so minimality is orientation-independent.
\end{solution}

\begin{exercise}[Minimal implies nonpositive \(K\)]\label{exr:vii19-ped-07}
Prove a minimal surface in \(\mathbb R^3\) has \(K\le0\).
\end{exercise}
\begin{hint}
Use \(k_2=-k_1\).
\end{hint}
\begin{solution}
Since \(H=0\), \(k_1+k_2=0\), so \(K=k_1k_2=-k_1^2\le0\).
\end{solution}

\begin{exercise}[Minimal and \(K=0\)]\label{exr:vii19-ped-08}
Show that a connected minimal surface patch with \(K=0\) is planar.
\end{exercise}
\begin{hint}
Combine \(H=0\) and \(K=0\) to determine principal curvatures.
\end{hint}
\begin{solution}
The two principal curvatures have sum and product zero, hence both vanish.  The shape operator is zero, so the normal is constant and the connected patch lies in a plane.
\end{solution}

\begin{exercise}[First variation principle]\label{exr:vii19-ped-09}
State the variational meaning of \(H=0\) for compactly supported normal variations.
\end{exercise}
\begin{hint}
Recall the first variation of area.
\end{hint}
\begin{solution}
The first variation is proportional to \(-\int 2H\,\varphi\,dA\) (up to convention); it vanishes for every compactly supported normal speed \(\varphi\) exactly when \(H=0\).
\end{solution}

\subsection*{Counterexamples and hypothesis tests}

\begin{exercise}[Minimal means zero Gaussian curvature]\label{exr:vii19-ped-10}
Test with the catenoid.
\end{exercise}
\begin{hint}
Minimality gives opposite principal curvatures.
\end{hint}
\begin{solution}
False.  Nonplanar minimal surfaces typically have \(K<0\).
\end{solution}

\begin{exercise}[Every ruled surface is minimal]\label{exr:vii19-ped-11}
Test with a cylinder.
\end{exercise}
\begin{hint}
A cylinder is ruled.
\end{hint}
\begin{solution}
False.  The cylinder is ruled but has nonzero mean curvature.
\end{solution}

\begin{exercise}[Every harmonic graph is exactly minimal]\label{exr:vii19-ped-12}
Test beyond the small-slope regime.
\end{exercise}
\begin{hint}
Compare the nonlinear minimal equation with the Laplace equation.
\end{hint}
\begin{solution}
False in general.  Harmonicity is the linearized equation; the exact minimal graph equation contains nonlinear first-derivative factors.
\end{solution}

\subsection*{Applications and investigations}

\begin{exercise}[Soap films]\label{exr:vii19-ped-13}
Why do ideal soap films model minimal surfaces?
\end{exercise}
\begin{hint}
Surface tension drives area decrease subject to boundary constraints.
\end{hint}
\begin{solution}
At equilibrium the first variation of area vanishes, leading to zero mean curvature away from singularities and external forces.
\end{solution}

\begin{exercise}[Small-slope approximation]\label{exr:vii19-ped-14}
Why does the Laplace equation approximate minimal graphs with small gradients?
\end{exercise}
\begin{hint}
Expand the nonlinear coefficients for \(|\nabla f|\ll1\).
\end{hint}
\begin{solution}
Terms quadratic in \(f_x,f_y\) are negligible, leaving \(f_{xx}+f_{yy}\approx0\).
\end{solution}

\subsection*{Challenge problems}

\begin{exercise}[No compact boundaryless minimal surface in Euclidean space]\label{exr:vii19-ped-15}
Explain why there is no nonconstant compact minimal surface without boundary immersed in \(\mathbb R^3\).
\end{exercise}
\begin{hint}
On a minimal surface, coordinate functions are harmonic; use the maximum principle.
\end{hint}
\begin{solution}
Each ambient coordinate restricted to a compact boundaryless minimal surface is harmonic and hence constant, forcing the immersion to be constant, a contradiction.
\end{solution}

\begin{exercise}[Minimal principal-curvature symmetry]\label{exr:vii19-ped-16}
Show that at every nonplanar minimal point the principal directions carry equal and opposite normal curvatures.
\end{exercise}
\begin{hint}
Use \(H=0\) and nonzero shape operator.
\end{hint}
\begin{solution}
\(k_1+k_2=0\) gives \(k_2=-k_1\), and nonplanarity at the point means \(k_1\neq0\); the two principal bends therefore have equal magnitude and opposite sign.
\end{solution}

% END VOL07-EXPANSION VII19
